\begin{lemma}[Differentiability of $u$ at a maximum point of the relabelled difference]
Let $N\in\mathbb N$, let $\Omega\subseteq\mathbb R^{N}$, let $\lambda>0$ and let
$u,v:\mathbb R^{N}\to\mathbb R$. Write $\psi(s):=s+\lambda e^{s}$ and assume
\begin{enumerate}
\item $\psi\circ u$, i.e.\ $y\mapsto u(y)+\lambda e^{u(y)}$, is locally semi-convex on
  $\Omega$ (\noderef{LocallySemiConvexOn});
\item $-v$ is locally semi-convex on $\Omega$;
\item $x\in\Omega$ and $x$ maximizes $y\mapsto\big(u(y)+\lambda e^{u(y)}\big)-v(y)$ over
  $\Omega$.
\end{enumerate}
Then there is a continuous linear functional $D$ on $\mathbb R^{N}$ such that $u$ is
Fr\'echet differentiable at $x$ with derivative $D$.
\end{lemma}
\begin{proof}
\noindent\textbf{Step 1: the relabelling $\psi$.} The function $s\mapsto s$ has derivative
$1$ everywhere and $s\mapsto\lambda e^{s}$ has derivative $\lambda e^{s}$ (the derivative of
the real exponential is itself, multiplied by the constant $\lambda$), so $\psi$ is
differentiable on $\mathbb R$ with
\[
 \psi'(s)=1+\lambda e^{s}\qquad(s\in\mathbb R). \tag{1}
\]
This is an elementary computation for an explicit function; no separate node is needed.
Since $\lambda>0$ and $e^{s}>0$ we have $\psi'(s)\ge1$ for every $s$. Put
\[
 t:=1+\lambda e^{u(x)} ,
\]
so that $t\ge1$, hence $t>0$, $t\ne0$ and $t^{-1}>0$.

\medskip\noindent\textbf{Step 2: $\psi$ is expanding.} For all $s,s'\in\mathbb R$,
\[
 |s-s'|\ \le\ |\psi(s)-\psi(s')| . \tag{2}
\]
Indeed $\psi(s)-\psi(s')=(s-s')+\lambda\big(e^{s}-e^{s'}\big)$. If $s'\le s$ then
$s-s'\ge0$, and $e^{s'}\le e^{s}$ with $\lambda>0$ gives
$\lambda(e^{s}-e^{s'})\ge0$; so $\psi(s)-\psi(s')\ge0$ and
$|\psi(s)-\psi(s')|=(s-s')+\lambda(e^{s}-e^{s'})\ge s-s'=|s-s'|$. If $s\le s'$ both
summands are $\le0$ and the same computation with reversed signs gives
$|\psi(s)-\psi(s')|=-(s-s')-\lambda(e^{s}-e^{s'})\ge-(s-s')=|s-s'|$.

\medskip\noindent\textbf{Step 3: differentiability of $\psi\circ u$ at $x$.} The first two
hypotheses of the lemma say that $\psi\circ u$ and $-v$ are locally semi-convex on $\Omega$,
and the third says that $x\in\Omega$ lies in
$\mathrm{argmax}_\Omega\big(\psi\circ u-v\big)$. These are
exactly the hypotheses of part (2) of \noderef{SemiConvexDiffAtMax} for the pair
$(\psi\circ u,\,v)$ on $\Omega$, so that node supplies a continuous linear functional
$D_\ell$ with
\[
 \psi\circ u\ \text{Fr\'echet differentiable at }x\ \text{with derivative }D_\ell
 \tag{3}
\]
(the node also gives that $v$ has the same derivative at $x$, which we do not use). Put
$L:=\|D_\ell\|+1$, so $L\ge1>0$.

\medskip\noindent\textbf{Step 4: $t^{-1}D_\ell$ is the derivative of $u$ at $x$.} We verify
the Fr\'echet condition in its $o(\|y-x\|)$ form: for every $c>0$ there is $\delta>0$ with
\[
 \Big|u(y)-u(x)-\big(t^{-1}D_\ell\big)(y-x)\Big|\ \le\ c\|y-x\|
 \qquad\text{whenever }\|y-x\|<\delta. \tag{4}
\]
Fix $c>0$. From (3), written in the same $o$-form, choose
\begin{align}
 \delta_1>0\ \text{with}\quad
 \big|\psi(u(y))-\psi(u(x))-D_\ell(y-x)\big|&\le1\cdot\|y-x\|
 &&\text{for }\|y-x\|<\delta_1, \tag{5a}\\
 \delta_2>0\ \text{with}\quad
 \big|\psi(u(y))-\psi(u(x))-D_\ell(y-x)\big|&\le\tfrac{c\,t}2\|y-x\|
 &&\text{for }\|y-x\|<\delta_2, \tag{5b}
\end{align}
the accuracies $1$ and $\frac{c\,t}2$ being positive. From Step 1, $\psi$ has derivative $t$
at the point $u(x)$ by (1); writing that one-dimensional statement in $o$-form with the
positive accuracy $\frac{c\,t}{2L}$, choose $\eta>0$ with
\[
 \Big|\psi(s)-\psi(u(x))-\big(s-u(x)\big)t\Big|\le\frac{c\,t}{2L}\,\big|s-u(x)\big|
 \qquad\text{for }|s-u(x)|<\eta. \tag{6}
\]
Put
\[
 \delta:=\min\Big\{\delta_1,\ \delta_2,\ \frac{\eta}{L+1}\Big\}>0
\]
and let $y$ satisfy $\|y-x\|<\delta$; then (5a), (5b) apply and $\|y-x\|<\eta/(L+1)$.

\emph{(4a) The crude Lipschitz bound.} By (2) applied to $u(y)$ and $u(x)$, then the
triangle inequality and (5a) together with $|D_\ell(y-x)|\le\|D_\ell\|\,\|y-x\|$,
\[
 |u(y)-u(x)|\le\big|\psi(u(y))-\psi(u(x))\big|
  \le\big|\psi(u(y))-\psi(u(x))-D_\ell(y-x)\big|+\big|D_\ell(y-x)\big|
\]
\[
  \le\|y-x\|+\|D_\ell\|\,\|y-x\|=L\|y-x\| . \tag{7}
\]

\emph{(4b) $u(y)$ lies in the $\eta$-ball about $u(x)$.} Multiplying
$\|y-x\|<\eta/(L+1)$ by $L>0$ and using (7),
\[
 |u(y)-u(x)|\le L\|y-x\|\le L\cdot\frac{\eta}{L+1}<\eta ,
\]
the last inequality because $L\eta<(L+1)\eta$ (as $\eta>0$) and $L+1>0$.

\emph{(4c) The relabelling estimate at $s:=u(y)$.} By (4b) we may use (6) at $s=u(y)$, and
then bound $|u(y)-u(x)|$ by (7), the factor $\frac{c\,t}{2L}$ being non-negative:
\[
 \Big|\psi(u(y))-\psi(u(x))-\big(u(y)-u(x)\big)t\Big|
  \le\frac{c\,t}{2L}\big|u(y)-u(x)\big|
  \le\frac{c\,t}{2L}\cdot L\|y-x\|=\frac{c\,t}2\|y-x\| . \tag{8}
\]

\emph{(4d) Combination.} The algebraic identity
\[
 \big(u(y)-u(x)\big)t-D_\ell(y-x)
  =-\Big(\psi(u(y))-\psi(u(x))-\big(u(y)-u(x)\big)t\Big)
   +\Big(\psi(u(y))-\psi(u(x))-D_\ell(y-x)\Big)
\]
holds because the two copies of $\psi(u(y))-\psi(u(x))$ cancel. By the triangle inequality,
$|-w|=|w|$, and then (8) and (5b),
\[
 \Big|\big(u(y)-u(x)\big)t-D_\ell(y-x)\Big|
  \le\frac{c\,t}2\|y-x\|+\frac{c\,t}2\|y-x\|=c\,t\,\|y-x\| . \tag{9}
\]

\emph{(4e) Division by $t$.} Since $t\ne0$, and
$\big(t^{-1}D_\ell\big)(y-x)=t^{-1}D_\ell(y-x)$ by homogeneity of $D_\ell$,
\[
 u(y)-u(x)-\big(t^{-1}D_\ell\big)(y-x)
  =t^{-1}\Big(\big(u(y)-u(x)\big)t-D_\ell(y-x)\Big).
\]
Taking absolute values, using $t^{-1}>0$ and then (9),
\[
 \Big|u(y)-u(x)-\big(t^{-1}D_\ell\big)(y-x)\Big|
  =t^{-1}\Big|\big(u(y)-u(x)\big)t-D_\ell(y-x)\Big|
  \le t^{-1}\cdot c\,t\,\|y-x\|=c\|y-x\| ,
\]
which is (4) with this $\delta$.

Since $c>0$ was arbitrary, $u$ is Fr\'echet differentiable at $x$ with derivative
$D:=t^{-1}D_\ell$, which is the assertion.
\end{proof}
